\section{随机微分方程（SDE）基础}

\subsection{SDE 的一般形式}

在连续时间域中，扩散模型的前向过程可以表示为如下 SDE：

$$
\mathrm{d}\mathbf{x} = f(\mathbf{x}, t)\,\mathrm{d}t + g(t)\,\mathrm{d}\mathbf{w}
$$

其中各符号含义如下：

\begin{itemize}
    \item $\mathbf{x}$ —— 数据变量（如图像），是关于时间 $t$ 的随机过程；
    \item $f(\mathbf{x}, t)$ —— \textbf{漂移系数}（drift coefficient），决定了 $\mathbf{x}_t$ 的均值随时间的演化方向；
    \item $g(t)$ —— \textbf{扩散系数}（diffusion coefficient），控制注入随机性的强度；
    \item $\mathbf{w}$ —— \textbf{标准 Wiener 过程}（维纳过程），又称布朗运动（Brownian motion）；
    \item $\mathrm{d}\mathbf{w}$ —— Wiener 过程的微小增量。
\end{itemize}

\begin{importantbox}{漂移与扩散的直觉}
\textbf{漂移项} $f(\mathbf{x}, t)\,\mathrm{d}t$ 可以理解为一个确定性的"力"，它将每个数据点沿某个方向移动（例如向原点收缩）。\textbf{扩散项} $g(t)\,\mathrm{d}\mathbf{w}$ 则注入随机噪声，使得 $\mathbf{x}_t$ 的轨迹变得锯齿状（jagged）。两者结合，就产生了从数据分布到噪声分布的连续映射。
\end{importantbox}

\subsection{Wiener 过程（布朗运动）}

Wiener 过程 $\mathbf{w}(t)$ 是 SDE 中随机性的来源，它具有以下四个特征性质：

\begin{enumerate}
    \item \textbf{初始条件}：$\mathbf{w}(0) = 0$（几乎必然，即 $P(\mathbf{w}(0)=0) = 1$）；
    \item \textbf{增量的正态性}：对任意 $0 \le s < t$，增量 $\mathbf{w}(t) - \mathbf{w}(s) \sim \mathcal{N}(0, (t-s)\mathbf{I})$；
    \item \textbf{独立增量}：不重叠时间区间上的增量相互独立；
    \item \textbf{路径连续性}：$\mathbf{w}(t)$ 关于 $t$ 几乎必然连续。
\end{enumerate}

\begin{knowledgebox}{Wiener 过程本身就是一个 SDE}
如果令 $f = 0$、$g = 1$，则 SDE 退化为 $\mathrm{d}\mathbf{x} = \mathrm{d}\mathbf{w}$，其解就是 Wiener 过程本身。因此，布朗运动可以看作最简单的 SDE——没有漂移，只有纯粹的随机扩散。
\end{knowledgebox}

\subsection{Taylor 展开与近似}

在将离散前向过程转化为连续 SDE 时，Taylor 展开是核心数学工具。给定光滑函数 $f$，其一阶近似为：

$$
f(t + \Delta t) \approx f(t) + f'(t) \cdot \Delta t
$$

例如，对于 $f(\epsilon) = \sqrt{1 - \epsilon}$，当 $\epsilon$ 很小时：

$$
\sqrt{1 - \epsilon} \approx 1 - \frac{\epsilon}{2}
$$

这一近似在推导 VP-SDE 的 $f(t)$ 和 $g(t)$ 时被反复使用。当时间步长 $\Delta t \to 0$，高阶项消失，从而得到精确的微分方程形式。

\begin{warningbox}{连续近似的前提}
从离散到连续的转化依赖于时间步长足够小（$\Delta t \to 0$），噪声调度函数（如 $\beta_t$）足够光滑的假设。对于实际实现中的有限步数离散化，连续 SDE 只是一个近似——步数越多，近似越精确。这也解释了为什么增加离散步数通常能改善生成质量。
\end{warningbox}

\subsection{本章小结}

SDE 通过漂移项和扩散项描述了数据在连续时间下的随机演化。Wiener 过程为其提供了标准的随机驱动。通过 Taylor 展开，我们可以从离散 Markov 链的转移概率出发，推导出对应的连续 SDE。


